%% Responsable: Dylber Cabrera
\section{Algoritmos aproximados}
\label{sec:aproximados}

El aproximado principal es la búsqueda por coordenadas (D3). Es la técnica que se
implementó y se ejecutó con el modelo real; sus resultados están en la
sección~\ref{sec:experimentos}. El reparto uniforme frente a UCB1 no se desarrolla en esta
entrega (sección~\ref{sec:reparto}) y aquí no se le atribuye ninguna cota.

\subsection{Descripción del algoritmo}

Se usa la notación de la sección~\ref{sec:formalizacion}: $L$ ranuras, la ranura $i$ con
$k_i$ opciones y $\Acc(c)$ la fracción de instancias del lote $I$ que $c$ cumple por
completo. Para cada ranura $i$ hay un conjunto de opciones permitidas
$O_i\subseteq\{0,\dots,k_i-1\}$, que por defecto es todo el rango. Reducirlo es la forma de
aplicar lo que enseñan las trazas de fallos (sección~\ref{sec:experimentos}). Sea $\pi$ un
orden de las ranuras y $P$ el tope de pasadas. Escribimos $c[i{:=}o]$ para la configuración
que sale de $c$ cambiando solo la ranura $i$ a la opción $o$.

\begin{algorithm}[h]
\caption{Coordenadas con un arranque}
\label{alg:coordenadas}
\begin{algorithmic}[1]
\Require orden $\pi$, opciones permitidas $O_i$, tope de pasadas $P$, presupuesto $B$
\State elegir $c$ al azar, con $c_i\in O_i$ para toda ranura $i$ \Comment{semilla fija}
\State $a\gets\Acc(c)$
\For{$t=1,\dots,P$}
  \State $\textit{mejoró}\gets\text{falso}$
  \For{cada ranura $i$ en el orden $\pi$}
    \For{cada opción $o\in O_i$ con $o\neq c_i$}
      \State $c'\gets c[i{:=}o]$;\quad $a'\gets\Acc(c')$
      \If{$a'>a$} \Comment{mejora estricta}
        \State $c\gets c'$;\quad $a\gets a'$;\quad $\textit{mejoró}\gets\text{verdadero}$
      \EndIf
    \EndFor
  \EndFor
  \If{\textbf{no} \textit{mejoró}} \State \textbf{salir} \Comment{pasada completa sin cambios} \EndIf
\EndFor
\State \Return $c$
\end{algorithmic}
\end{algorithm}

Dos detalles de la implementación (\texttt{coordenadas.py}) que el pseudocódigo respeta. La
comparación $o\neq c_i$ usa la opción \emph{actual} en ese momento: tras aceptar un cambio
dentro de una ranura, la opción anterior vuelve a probarse más adelante en esa misma
ranura (sale del caché, sin generaciones nuevas). Y si el presupuesto $B$ se agota, el
algoritmo devuelve la mejor configuración vista hasta ahí, que es la actual por el
lema~\ref{lem:monotonia}. Varios arranques son repeticiones independientes del algoritmo
desde puntos de partida distintos; se devuelve la mejor de las configuraciones finales.

\subsection{Propiedades}

\begin{lema}[Monotonía]\label{lem:monotonia}
La precisión de la configuración actual no decrece, y cada cambio la aumenta de forma
estricta. En particular, la configuración devuelta es al menos tan buena como la de partida.
\end{lema}
\begin{proof}
La configuración solo cambia en la línea con la condición $a'>a$, y en ese caso $a$ pasa a
$a'$, estrictamente mayor. Si no hay cambio, $c$ y $a$ quedan igual.
\end{proof}

\begin{lema}[Terminación]\label{lem:terminacion}
Con $\Acc$ determinista, el algoritmo hace a lo sumo $n$ cambios y a lo sumo
$\min(P,\,n+1)$ pasadas.
\end{lema}
\begin{proof}
$\Acc(c)$ toma valores en $\{0,\tfrac1n,\dots,1\}$, así que cada cambio lo sube al menos
$\tfrac1n$, y como parte de al menos $0$ y no pasa de $1$, solo cabe $n$ veces. Toda pasada
que no es la última contiene al menos un cambio, de modo que hay a lo sumo $n$ pasadas con
cambios más la pasada final sin cambios.
\end{proof}

\begin{teorema}[Óptimo local por ranura]\label{teo:optimo-local}
Si el algoritmo termina porque una pasada completa no cambió nada, la configuración $c$
devuelta cumple $\Acc(c[i{:=}o])\le\Acc(c)$ para toda ranura $i$ y toda opción $o\in O_i$.
\end{teorema}
\begin{proof}
En esa pasada la configuración no cambió, así que todos los candidatos $c[i{:=}o]$ se
compararon contra el mismo $a=\Acc(c)$ y ninguno lo superó estrictamente. Cada candidato con
$o\in O_i$ y $o\ne c_i$ se evaluó en esa pasada; el caso $o=c_i$ da $c$ mismo.
\end{proof}

El teorema depende de que $\Acc$ devuelva el mismo valor cada vez que se evalúa una
configuración. Con un oráculo exacto es inmediato; con el modelo real se apoya en la
temperatura $0{,}0$ y en el caché, que devuelve lo ya generado (D2).

Es un óptimo \emph{por ranura}: ningún cambio de una sola ranura mejora, pero no se prueba
ningún cambio simultáneo de dos ranuras. Además, si $O_i$ se reduce a un subconjunto, el
óptimo es local respecto a ese subespacio, no respecto a todas las opciones. La
sección~\ref{sec:experimentos} muestra que esto ocurrió en el arranque 1.

\subsection{Costo $T(n)$}

El costo se mide en generaciones, que es lo que cobra el presupuesto $B$ (D11): una
generación tarda unos $1{,}99$ s en la GPU del proyecto, mucho más que cualquier otra
operación del algoritmo. Evaluar una configuración sobre $I$ cuesta $n$ generaciones, una
por instancia.

\begin{teorema}[Costo]\label{teo:costo}
Sea $p\le\min(P,n+1)$ el número de pasadas realizadas. El número de evaluaciones de
configuraciones es a lo sumo $1+p\sum_i|O_i|$, y las generaciones, a lo sumo
\[
  T(n)\;\le\; n\Bigl(1+p\sum_{i=1}^{L}|O_i|\Bigr)\;=\;O(p\cdot L\cdot k_{\max}\cdot n).
\]
\end{teorema}
\begin{proof}
En una pasada, el recorrido de la ranura $i$ examina cada opción de $O_i$ una vez, y
evalúa todas menos la actual, así que hace a lo sumo $|O_i|$ evaluaciones. Sumar sobre las
ranuras da $\sum_i|O_i|$ por pasada, más la evaluación inicial. Cada evaluación son $n$
generaciones. El caché solo puede bajar la cifra, porque no cobra una pareja
(configuración, instancia) ya generada.
\end{proof}

Con $P$ constante, el costo es lineal en $n$. Sin tope, el lema~\ref{lem:terminacion} da
$p\le n+1$ y el peor caso sube a $O(L\cdot k_{\max}\cdot n^2)$. El costo no depende de
$|\conf|=\prod_i k_i$, a diferencia del recorrido exhaustivo, que es $O(|\conf|\cdot n)$.

Para el caso del proyecto, $L=5$, $k_i=8$ y $n=100$:
\begin{center}
\begin{tabular}{lr}
\toprule
Concepto & Generaciones, cota \\
\midrule
Una pasada, sin exclusiones: $5\cdot 8\cdot 100$ & $4\,000$ \\
Un arranque con $P=3$: $100\,(1+3\cdot 40)$ & $12\,100$ \\
Presupuesto $B$ (D17) & $14\,400$ \\
\bottomrule
\end{tabular}
\end{center}
Dos arranques de $P=3$ podrían sumar, en el peor caso, más que $B$. Lo que acota el gasto
real es el contador de generaciones, que lanza \texttt{PresupuestoAgotado} al llegar a $B$.
Lo observado fue menor y está en la sección~\ref{sec:experimentos}: el arranque 0 evaluó
$92$ configuraciones distintas ($9\,200$ generaciones) y el arranque 1, $28$ ($2\,800$).

\subsection{Cota de calidad: no hay razón de aproximación}

El algoritmo garantiza el lema~\ref{lem:monotonia} y el teorema~\ref{teo:optimo-local},
pero no una razón de aproximación respecto al óptimo.

\begin{teorema}[Sin razón de aproximación]\label{teo:contraejemplo}
No existe una constante $\alpha>0$ tal que el algoritmo~\ref{alg:coordenadas} devuelva
siempre $c$ con $\Acc(c)\ge\alpha\,\Acc(c^*)$, donde $c^*$ es un óptimo.
\end{teorema}
\begin{proof}
Sean $L=2$ ranuras con dos opciones cada una, y un lote de $n$ instancias idénticas en el
que cumplir la instancia exige dos condiciones: una la satisface solo la opción $1$ de la
ranura 1 y la otra, solo la opción $1$ de la ranura 2. Por el puntaje todo-o-nada,
$\Acc(1,1)=1$ y $\Acc(c)=0$ para las otras tres configuraciones. Si el punto de partida es
$(0,0)$, sus dos vecinos tienen precisión $0$, no hay mejora estricta y el algoritmo termina
en $(0,0)$ con $\Acc=0$, frente a $\Acc(c^*)=1$. La cantidad de ranuras y opciones del
proyecto solo agrega ranuras irrelevantes y no altera este argumento.
\end{proof}

Este es un ejemplo construido en el modelo, no una observación con el modelo real. Lo que
lo hace posible es el puntaje todo-o-nada, que crea mesetas donde ningún vecino mejora. Lo
que sí se observó con el modelo real va en la misma dirección: dos arranques con puntos de
partida y espacios distintos terminaron en 20\,\% y 17\,\% (sección~\ref{sec:experimentos}),
y la mejor opción de estilo depende de la estrategia elegida. Por eso el riesgo de quedar
en un óptimo local se mide con varios arranques y no se declara cubierto.

\subsection{Qué no se afirma}

\begin{itemize}
  \item No se afirma que la configuración encontrada sea un óptimo global ni del
  subcatálogo. Es la mejor encontrada.
  \item No hay garantía probabilística de haber hallado el óptimo: el análisis de muestreo
  de la sección~\ref{sec:reduccion} muestra que una garantía uniforme sobre todo el espacio
  es inviable con este presupuesto.
  \item La precisión de búsqueda es la de un lote fijo de $100$ instancias y no se extiende
  sin supuestos a otras instancias.
\end{itemize}
